\section{Method}

\subsection{The classifier}

Let the training set have positive examples $x_i$ with $i \in P$ ($n_+$ of them) and negative examples with $i \in M$ ($n_-$ of them), and let $K$ be the kernel. Each positive example gets weight $\alpha_i = 1/n_+$, each negative example gets weight $\alpha_i = -1/n_-$, and the classifier is
\begin{equation}
h(x) = \operatorname{sign}\Bigl(\sum_i \alpha_i K(x_i, x) - b\Bigr), \label{eq:rule}
\end{equation}
with ties labeled $+1$ and the bias
\begin{equation}
b = \frac{1}{2}\Bigl(\frac{1}{n_+^2}\sum_{i,j \in P} K(x_i, x_j) - \frac{1}{n_-^2}\sum_{i,j \in M} K(x_i, x_j)\Bigr). \label{eq:bias}
\end{equation}

Let $\mu_+$ and $\mu_-$ be the means of the positive and negative training points in the kernel's feature space. The quantity inside the sign in \eqref{eq:rule} equals $\tfrac{1}{2}(\lVert\phi(x)-\mu_-\rVert^2 - \lVert\phi(x)-\mu_+\rVert^2)$, so the rule assigns $x$ to the class whose mean is nearer. This is the simple kernel classifier described by Sch\"olkopf and Smola \cite[Sec.~1.2]{scholkopf2002}.

Three cases follow. With the dot product, the rule is the nearest-centroid rule on the raw measurements. With a polynomial kernel, the decision boundary is a fixed polynomial of degree 2 or 3 determined by moments of the two classes. With the Gaussian kernel, each class sum is proportional to a Gaussian Parzen density estimate, so the rule compares the two classes' density estimates, each weighted equally, and thresholds their difference at a value set by $b$. That is the plug-in Bayes rule for equal class priors only when $b = 0$. The classifier outputs labels, not probabilities.

\subsection{Kernels}

The four kernels are
\begin{align}
K_{\mathrm{dot}}(x, y) &= x \cdot y, \nonumber\\
K_{2}(x, y) &= (x \cdot y + 1)^2, \nonumber\\
K_{3}(x, y) &= (x \cdot y + 1)^3, \nonumber\\
K_{\mathrm{G}}(x, y) &= \exp\bigl(-\lVert x - y\rVert^2 / 2\sigma^2\bigr), \label{eq:kernels}
\end{align}
with $\sigma = 1/\sqrt{2} \approx 0.71$ for the Gaussian. The width was set by the assignment and was the same for every run. It was not tuned, so no choice was made on test data. Problem 1 uses the dot product only.

\subsection{Data and problems}

The data are the 150 Iris flowers distributed with scikit-learn \cite{pedregosa2011}: 50 each of setosa, versicolor, and virginica, with four measurements in centimeters (sepal length, sepal width, petal length, and petal width), used unscaled. Four binary problems are built from them. Problem 1 is setosa against the other two species. Problem 2 is virginica against the other two, and Problem 3 is versicolor against the other two. Problem 4 is versicolor against virginica, with setosa removed (\N{p4_samples} flowers). Setosa is linearly separable from the other two species, and versicolor and virginica are not linearly separable from each other, so Problem 1 is the only linearly separable problem. Problems 1 to 3 pit one species against the other two, so their classes have a 1:2 ratio.

\subsection{Evaluation}

The training share takes five values: 90\%, 80\%, 70\%, 60\%, and 50\%. For each share, \N{n_splits_ratio} random splits are drawn with scikit-learn's \texttt{train\_test\_split}, using seeds 0 to \N{seed_last} and no stratification, as in the original study. Every kernel and problem uses the same splits. The 13 problem--kernel pairs, five shares, and \N{n_splits_ratio} splits make \N{n_fits} fits. The results report the mean test accuracy over the splits, a 95\% confidence interval (CI) for that mean, and the 2.5th to 97.5th percentiles of single splits. The interval is the mean plus or minus $1.96\,\mathrm{SD}/\sqrt{n}$, where $n$ is \N{n_splits_ratio}. At one share, SD is the standard deviation over the splits. When the shares are pooled, the five splits drawn with one seed are not independent (they share a shuffle, and their test sets are nested), so SD is the standard deviation of the \N{n_splits_ratio} per-seed means, each the mean of one seed's five test accuracies. All splits resample the same 150 flowers, so the interval describes the choice of splits only, not new flowers.

The baseline is the majority-class rate: the larger class's share of each test set, averaged over splits. It is at least what always predicting one fixed label would score on each split.

The original study used seed 42 at each share, which is one of the \N{n_splits_ratio} seeds. Its five splits came from one shuffle, so each test set sits inside the next larger one. Its \N{n_runs} runs (13 pairs at five shares) are kept as a comparison: the ``Original'' column of Table~\ref{tab:kernels} and the open markers in the figures. The evaluation code checks that its classifier gives the original code's predictions and accuracies on all \N{n_runs} original runs.

As an added comparison, the same splits are evaluated with each feature standardized to mean 0 and standard deviation 1, using the mean and standard deviation of the problem's own training rows.
